% Generated draft: requires mathematical and source review.
\documentclass[11pt]{article}
\usepackage[margin=25mm]{geometry}
\usepackage{fontspec}
\setmainfont{Latin Modern Roman}
\usepackage{amsmath,amssymb,mathtools,unicode-math}
\setmathfont{Latin Modern Math}
\usepackage{enumitem,needspace,xurl,hyperref}
\setlength{\parindent}{0pt}
\setlength{\parskip}{5pt}
\setlength{\emergencystretch}{3em}
\newcommand{\N}{\mathbb{N}}
\newcommand{\Z}{\mathbb{Z}}
\newcommand{\Q}{\mathbb{Q}}
\newcommand{\R}{\mathbb{R}}
\newcommand{\C}{\mathbb{C}}
\newcommand{\F}{\mathbb{F}}
\newcommand{\A}{\mathbb{A}}
\newcommand{\PP}{\mathbb{P}}
\newcommand{\E}{\mathbb{E}}
\newcommand{\eps}{\varepsilon}
\newcommand{\dd}{\,\mathrm{d}}
\begin{document}
\section*{UnsolvedMath ID: 30003455}
\textbf{Source problem code:} catalogue entry
\subsection{Definitions and mathematical statement}
The supplied record does not define the univariate smoothness classes to be generalized. A smoothness class here refers informally to a collection of functions subject to specified regularity conditions; those conditions, the domain, the range, and any norm or seminorm are not specified. A multivariate analogue would concern functions of several variables, but neither the number of variables nor the meaning of a natural analogue is fixed.

\textbf{Research question.} What is a natural multivariate analogue of univariate smoothness classes, and how should functions belonging to those classes be approximated?

A precise quantified approximation problem cannot be extracted from the supplied statement alone. It does not specify the underlying univariate classes, multivariate domains, regularity parameters, boundary conditions, admissible approximants, approximation error, or approximation budget. The accompanying background mentions optimal approximation, mixed smoothness, and associated \(n\)-width questions, but supplies no definitions or particular extremal problem. These references provide context and do not determine additional hypotheses or an optimality criterion. Consequently, this draft preserves an open-ended research question rather than asserting a specific theorem or conjectured approximation rate.
\subsection{Short English statement}
Identify natural smoothness classes for functions of several variables and determine suitable ways to approximate their members. The supplied record leaves the classes, domains, error criteria, and meaning of optimality unspecified, so it does not determine a single precise conjecture.
\subsection{Sources}
\begin{itemize}
\item Local UnsolvedMath registry, permanent ID \(30003455\), title: Multivariate Smoothness Classes and Their Approximation; exact source problem code: \texttt{catalogue entry}. Collection: Oberwolfach Reports Open Problems. No individual proposer is identified; the supplied proposed year is 2017.
\item Supplied report attribution: \emph{Multiscale and High-Dimensional Problems} (2018), Oberwolfach Reports, pp.\ 1001--1051, DOI 10.4171/owr/2017/17. The source URL and the literature source labelled ``Original Oberwolfach report'' are \url{https://doi.org/10.4171/owr/2017/17}. The supplied background describes Oberwolfach Reports as documenting mathematical workshops and their open-problem discussions. The report contents have not been independently consulted for this draft.
\item Supplied literature assessment, checked on exactly 2026-08-21: the record labels the problem open and describes natural multivariate smoothness classes, optimal approximation, mixed-smoothness approximation, and associated \(n\)-width questions as an active research area rather than a single settled theorem. This is supplied dated curation, not independently verified literature status. The supplied background expressly calls for independent verification of time-sensitive status and cited claims.
\end{itemize}
\end{document}
